\chapter{书系中的位置}
\label{app:series-roadmap}

本册将原第二卷中前二〇六年至二二〇年的材料拆出，作为能够独立阅读的汉帝国卷。它以楚汉战争结束后的战后秩序为入口，以汉魏禅代为出口；前二〇九至前二〇七年的秦末起事作为开国条件留在开篇背景中，但不改变本卷的年代边界。

下一册将从二二〇年的魏、蜀、吴并立条件继续，处理三国、西晋、东晋、十六国与南北朝。二二〇年被放在本册末尾，是因为禅代既结束东汉皇统，也把汉帝国留下的都城、尚书、军政、地方化和正统问题交给后续各政权。这样切分不把魏初看作与汉无关的全新开端，也避免让后续数百年的分裂史压缩本册对汉代制度、社会与边疆的讨论。

本卷把来源说明、事件索引和图件放在各自的阅读入口中：它们帮助读者从一段叙事追问证据和空间，而不要求读者进入制作过程。三国以后的内容没有被搬入本册；下一册才从二二〇年的断点继续展开。
